\chapter{Absolute quantum reduction and higher products}
\label{chap:vol2-absolute-quantum-reduction}

The relative reduction retains relations among the moment
equations. The absolute reduction must also retain the gauge
directions. Even a trivial coefficient module for
$\mathfrak{so}_3$ has a degree-three Lie cohomology class.
Multiplying this class by the relative degree-minus-one relation
class produces a degree-two class. Both disappear if one keeps
only invariant coefficients and discards the positive ghosts.

We compute the full polynomial traceless rank-two complex over
$R=\mathbb C[\hbar]$. Its cohomology is supported in degrees
$-1,0,2,3$. A contraction constructs this answer across the
discriminant and specifies every higher product.
Its projection cannot be a unital graded algebra map to the
cohomology ring. That obstruction has a fixed direction:
whether a zigzag of dg-algebra quasi-isomorphisms exists is a
separate question.

\section{The positive ghosts}
\label{sec:vol2-absolute-ghosts}

Use the traceless coordinates $x_i,y_i$ and antighosts $\psi_i$
of Chapter~\ref{chap:vol2-relative-quantum-reduction}, for
$1\leq i\leq3$. Their degrees are $0,0,-1$.
Adjoin positive ghosts $\gamma^i$ of degree $1$ and impose
\begin{equation}\label{eq:vol2-absolute-ghost-relations}
 [x_i,y_j]=\hbar\delta_{ij},\qquad
 \psi_i\gamma^j+\gamma^j\psi_i=\hbar\delta_i^j,
 \qquad
 \gamma^i\gamma^j=-\gamma^j\gamma^i,\quad
 \psi_i\psi_j=-\psi_j\psi_i.
\end{equation}
The even coordinates commute with the ghosts.
The $x_i$ commute among themselves, as do the $y_i$.
The odd squares are zero. Write
$\mathrm{Cl}_\hbar(\gamma,\psi)$ for this graded Clifford
algebra and put
\[
 \mathcal B_\hbar
 =W_\hbar(x,y)\otimes_R\mathrm{Cl}_\hbar(\gamma,\psi).
\]
Use symmetric Weyl symbols in $x,y$ and place all positive
ghosts to the left of all antighosts. The ordered products
$\gamma^I\psi_J$, for increasing subsets $I,J\subset\{1,2,3\}$,
form an $R$-basis of the ghost algebra.
Indeed the relations span it by these products. After inverting
$\hbar$, their action on $\Lambda(\gamma)$ by multiplication
and $\psi_i=\hbar\partial_{\gamma^i}$ is the full matrix
algebra, whose dimension is $2^6$. Thus they are independent
already over $R$.

For the rotation moment define
\begin{equation}\label{eq:vol2-absolute-brst-charge}
 \mu_i=\sum_{j,k}\epsilon_{ijk}x_jy_k,\qquad
 \Theta=\sum_i\gamma^i\mu_i
       -\frac12\sum_{i,j,k}\epsilon_{ijk}
                    \gamma^i\gamma^j\psi_k,
 \qquad Q=\frac{[\Theta,-]_{\mathrm{super}}}{\hbar}.
\end{equation}
Here $\epsilon_{123}=1$ is alternating, and the bracket is the
graded commutator for the full product.
Division means the polynomial quotient of a divisible numerator.
The algebra modulo $\hbar$ is graded commutative, so every such
commutator is divisible. The Weyl--Clifford basis makes the
quotient unique.

The moment relation is
$[\mu_i,\mu_j]=\hbar\sum_k\epsilon_{ijk}\mu_k$.
Together with Jacobi it gives $\Theta^2=0$, as in
Proposition~\ref{prop:vol2-integral-quantum-brst}.
One can also represent the ghosts on $\Lambda(\gamma)$ after
inverting $\hbar$. The charge becomes the Chevalley operator
for the representation $\mu_i/\hbar$, multiplied by $\hbar$.
Its square vanishes. Faithfulness and polynomial torsion-freeness
give the same identity before inversion.
Consequently $Q^2=0$, and $Q$ is a degree-one $R$-linear
derivation. These statements remain valid at $\hbar=0$ by
specializing the polynomial formulas.

The distinction from invariants already appears in the scalar
Chevalley complex. Its differential is
\begin{equation}\label{eq:vol2-scalar-ghost-differential}
 d\gamma^1=-\gamma^2\gamma^3,\qquad
 d\gamma^2=\gamma^1\gamma^3,\qquad
 d\gamma^3=-\gamma^1\gamma^2.
\end{equation}
It is an isomorphism from exterior degree one to degree two
and is zero in every other degree. Thus its cohomology is
$\Lambda(\chi)$, where
\[
 \chi=\gamma^1\gamma^2\gamma^3,\qquad |\chi|=3.
\]
For the full complex, the positive ghosts also act on the
polynomial coefficients. We next contract that action without
removing the scalar and top-ghost classes.

\section{Normal ordering and the two differentials}
\label{sec:vol2-absolute-normal-order}

Let
$K_{\mathrm{all}}=W_\hbar(x,y)\otimes\Lambda(\psi)$
be the space of coefficients with no positive ghost.
For the basis $e_i$ of $\mathfrak{so}_3$, use
$[e_i,e_j]=\sum_k\epsilon_{ijk}e_k$.
The simultaneous coefficient action $\rho_i$ is
\begin{equation}\label{eq:vol2-absolute-coefficient-action}
 \rho_i x_j=\sum_k\epsilon_{ijk}x_k,
 \qquad \rho_i y_j=\sum_k\epsilon_{ijk}y_k,
 \qquad \rho_i\psi_j=\sum_k\epsilon_{ijk}\psi_k.
\end{equation}
It extends by the derivation rule. On Weyl coefficients it is
$[\mu_i,-]/\hbar$. It does not act on $\gamma$ in this notation.
Each polynomial lies in a finite-dimensional rational rotation
module, so the coefficient representation is locally finite.

For $X\in\mathfrak{so}_3$, let $\mu_X$ and $\psi_X$ denote
the corresponding linear combinations. Define a degree-one
operator on the antighost complex by
\begin{align}\label{eq:vol2-absolute-chain}
 \partial_N(a\psi_{X_1}\cdots\psi_{X_q})
 ={}&\sum_{i=1}^q(-1)^{i-1}(a\star\mu_{X_i})
     \psi_{X_1}\cdots\widehat{\psi_{X_i}}\cdots\psi_{X_q}
       \notag\\
 &+\hbar\sum_{i<j}(-1)^{i+j}a\psi_{[X_i,X_j]}
     \psi_{X_1}\cdots\widehat{\psi_{X_i}}\cdots
                       \widehat{\psi_{X_j}}\cdots\psi_{X_q}.
\end{align}
The bracket factor in the second line stands first.
For example,
\[
 \partial_N(a\psi_i\psi_j)
 =a\star\mu_i\,\psi_j-a\star\mu_j\,\psi_i
          -\hbar a\psi_{[e_i,e_j]}.
\]
Applying $\partial_N$ again gives the moment commutator minus
the same Lie bracket. It is zero. In higher antighost degree,
the terms involving two moments cancel in the same way;
the terms involving three indices cancel by Jacobi.
This proves $\partial_N^2=0$ in every degree.
The moment map and the Lie bracket are equivariant, so
$\partial_N$ commutes with each $\rho_i$.

Let $\iota_i=\partial_{\gamma^i}$ be the left exterior
derivative. Define
\begin{equation}\label{eq:vol2-absolute-chevalley}
 d_{\mathrm{CE}}=\sum_i\gamma^i\rho_i
       -\frac12\sum_{i,j,k}\epsilon_{ijk}
                     \gamma^i\gamma^j\iota_k.
\end{equation}
This is the Chevalley cochain differential with coefficients in
$K_{\mathrm{all}}$. The Lie representation identity and Jacobi
give $d_{\mathrm{CE}}^2=0$.
On positive-ghost degree $p$, put
$\partial^{\mathrm{tot}}=(-1)^p\partial_N$.
Equivariance gives
$d_{\mathrm{CE}}\partial^{\mathrm{tot}}
 +\partial^{\mathrm{tot}}d_{\mathrm{CE}}=0$.

\begin{lemma}\label{lem:vol2-absolute-normal-order}
Normal ordering identifies the underlying complex with
\begin{equation}\label{eq:vol2-absolute-bicomplex}
 (\mathcal B_\hbar,Q)
  =\bigl(\Lambda(\gamma)\otimes K_{\mathrm{all}},
          d_{\mathrm{CE}}+\partial^{\mathrm{tot}}\bigr).
\end{equation}
On invariant positive-ghost-free coefficients, $Q$ is the
relative differential $\delta_\hbar$ of
\eqref{eq:vol2-relative-quantum-differential}.
\end{lemma}

\begin{proof}
Commute each positive ghost produced by $[\Theta,-]$ to the
left. The terms retaining one new ghost give the simultaneous
coefficient action. A contraction with an existing positive
ghost gives the second term of
\eqref{eq:vol2-absolute-chevalley}.
Removing one antighost gives the moment term in
\eqref{eq:vol2-absolute-chain}. Contracting the cubic ghost
term against two antighosts gives its bracket term.
Moving this degree-one operator past $p$ positive ghosts
gives $(-1)^p$.

On generators, for example,
\[
 Q\gamma^1=-\gamma^2\gamma^3,
 \qquad Q\psi_1=\mu_1-\gamma^2\psi_3+\gamma^3\psi_2,
 \qquad Qx_j=\sum_{i,k}\gamma^i\epsilon_{ijk}x_k.
\]
The same formula holds for $Qy_j$.
The ordered computation above applies to an arbitrary product
and hence proves the identity on the whole basis.

For invariant coefficients, $\sum\gamma^i\rho_i$ vanishes.
Invariance combines the ordered moment and bracket terms into
\[
 \frac12\sum_i\left(
     \mu_i\star\partial_{\psi_i}F
       +\partial_{\psi_i}F\star\mu_i\right).
\]
The first Moyal derivatives cancel. The second derivatives of
the bilinear moments give
$\hbar^2\sum\epsilon_{ijk}
 \partial_{x_j}\partial_{y_k}\partial_{\psi_i}/4$,
and all higher derivatives vanish. This is the stated
relative differential.
\end{proof}

Equation~\eqref{eq:vol2-absolute-bicomplex} is an equality of
complexes after choosing an ordered basis.
Multiplication still uses
\eqref{eq:vol2-absolute-ghost-relations}.
Replacing it by the exterior tensor product would remove
terms needed for both the ring comparison and the higher products.

\section{Contracting the coefficient action}
\label{sec:vol2-absolute-casimir}

The scalar Chevalley complex should remain. To contract its
nontrivial coefficient representations, define
\begin{equation}\label{eq:vol2-absolute-casimir}
 \mathsf C=\sum_i\rho_i^2,
 \qquad K=\sum_i\rho_i\iota_i.
\end{equation}
Here $\mathsf C$ acts only on $K_{\mathrm{all}}$.
The diagonal rotation action on both coefficients and positive
ghosts has a different Casimir.

\begin{lemma}\label{lem:vol2-absolute-casimir}
The operators satisfy
\[
 d_{\mathrm{CE}}K+Kd_{\mathrm{CE}}=\mathsf C,
 \qquad
 \partial^{\mathrm{tot}}K+K\partial^{\mathrm{tot}}=0.
\]
Moreover $\mathsf C$ commutes with both differentials.
\end{lemma}

\begin{proof}
Use $\iota_j\gamma^i+\gamma^i\iota_j=\delta_j^i$.
The anticommutator of $K$ with $\sum\gamma^i\rho_i$ is
\[
 \mathsf C+\sum_{i,j,k}\epsilon_{ijk}
                   \gamma^i\rho_k\iota_j.
\]
Its anticommutator with the bracket part of
$d_{\mathrm{CE}}$ is the negative of the second summand.
Indeed commuting $\iota_\ell$ through
$\gamma^i\gamma^j$ gives
$\delta_\ell^i\gamma^j-\delta_\ell^j\gamma^i$;
antisymmetry then gives the displayed cancellation.

The coefficient Casimir is central for the rotation action,
so it commutes with $d_{\mathrm{CE}}$.
Equivariance of $\partial_N$ gives its commutation with
$\rho_i$ and $\mathsf C$.
Finally $K$ lowers positive-ghost degree by one, reversing
the totalization sign of $\partial_N$. This proves the
second anticommutator identity.
\end{proof}

On an irreducible spin-$\ell$ coefficient representation,
$\mathsf C=-\ell(\ell+1)$, for $\ell\geq0$.
These are integer spins because the representation is rational
for $SO(3)$.
To check the normalization, put
$E=\rho_1+i\rho_2$, $F=-\rho_1+i\rho_2$ and
$H_{\mathrm{wt}}=2i\rho_3$. They satisfy
$[H_{\mathrm{wt}},E]=2E$, $[H_{\mathrm{wt}},F]=-2F$,
$[E,F]=H_{\mathrm{wt}}$, and
\[
 \mathsf C=-\frac14H_{\mathrm{wt}}^2-\frac12(EF+FE).
\]
On a highest-weight vector of weight $2\ell$, this is
$-\ell^2-\ell$. Complete reducibility in each finite coefficient
module follows by averaging a Hermitian form over the compact
real rotation group and taking invariant orthogonal complements.
Thus the eigenvalue decomposition used below is an actual direct
sum on every finite input module.
For an input contained in spins at most $L$, define
\begin{equation}\label{eq:vol2-absolute-spectral-polynomials}
 P_L(X)=\prod_{\ell=1}^L
               \left(1+\frac{X}{\ell(\ell+1)}\right),
 \qquad
 G_L(X)=\frac{1-P_L(X)}{X}+P_L'(0)P_L(X).
\end{equation}
Both expressions are polynomials.
The second has value zero at $X=0$ and value $1/X$ at
every nonzero Casimir eigenvalue in this range.
Thus
\[
 P_0=P_L(\mathsf C),\qquad G=G_L(\mathsf C)
\]
are, respectively, projection to coefficient invariants and
the inverse Casimir on their complement, extended by zero.
Enlarging $L$ gives the same operators on the original input.
This constructs them on the union of finite coefficient modules.
They commute with equivariant maps, including $\partial_N$.
No power of $\hbar$ or of the discriminant is inverted.

Set $H_0=GK$. Then
\[
 QH_0+H_0Q=1-P_0,\qquad
 P_0H_0=H_0P_0=0.
\]
To impose the square-zero side condition, replace $H_0$ by
\begin{equation}\label{eq:vol2-absolute-normalized-casimir}
 H_1=H_0d_{\mathrm{CE}}H_0.
\end{equation}
On the complement of $P_0$, write $d=d_{\mathrm{CE}}$ and
$h=H_0$. The identity $dh+hd=1$ gives
$dh^2=h^2d$. Hence
\[
 d(hdh)+(hdh)d=dh+hd=1,
 \qquad (hdh)^2=h(dh^2d)h=0.
\]
The total chain differential anticommutes with $h,d,h$ in
succession, and therefore with $hdh$.
Consequently
\begin{equation}\label{eq:vol2-absolute-first-contraction}
 QH_1+H_1Q=1-P_0,
 \qquad H_1^2=H_1P_0=P_0H_1=0.
\end{equation}

The image of $P_0$ is
$\mathrm{CE}^*(\mathfrak{so}_3,\mathbb C)
 \otimes\mathcal K_\hbar$,
with the tensor differential. The next step contracts its
first factor explicitly, then uses the relative comparison.

\section{The absolute contraction}
\label{sec:vol2-absolute-contraction}

Let $j_\gamma:\Lambda(\chi)\to\Lambda(\gamma)$ include
the scalar and top-ghost monomials, and let $p_\gamma$ project
onto them. The homotopy for
\eqref{eq:vol2-scalar-ghost-differential} is
\begin{equation}\label{eq:vol2-absolute-scalar-homotopy}
 h_\gamma(\gamma^2\gamma^3)=-\gamma^1,
 \quad h_\gamma(\gamma^1\gamma^3)=\gamma^2,
 \quad h_\gamma(\gamma^1\gamma^2)=-\gamma^3,
\end{equation}
and zero on other exterior degrees. Direct substitution gives
$dh_\gamma+h_\gamma d=1-j_\gamma p_\gamma$ and the
normalized side conditions.

Write $(\rho_\hbar,i_{\mathrm{rel}},H_{\mathrm{rel}})$
for the relative contraction in
Section~\ref{sec:vol2-relative-higher-inverse}.
Its target is
\[
 \mathcal H_\hbar
  =\bigl(W_\hbar(s,t)\otimes\Lambda(\theta)\bigr)^{\mathbb Z/2},
 \qquad [s,t]=\hbar,\quad |\theta|=-1.
\]
The involution negates $s,t,\theta$ and fixes $\hbar$.
Choose its displayed section so that
$i_{\mathrm{rel}}(s\theta)=\xi=x\cdot\psi$,
$i_{\mathrm{rel}}(t\theta)=\eta=y\cdot\psi$ and
$H_{\mathrm{rel}}(\xi\star\eta)=r=\psi_1\psi_2\psi_3$.
Its formulas use finite classical weight splittings and the
terminating coefficient operator $T^{-1}$.

Put $Z=\Lambda(\chi)\otimes\mathcal H_\hbar$, initially
as a graded $R$-module with zero differential.
Compose the three contractions: first
\eqref{eq:vol2-absolute-first-contraction}, then
$(j_\gamma,p_\gamma,h_\gamma)$ tensored with the relative
complex, and finally the relative contraction tensored with
$\Lambda(\chi)$.
On the last tensor product, for $e\in\{0,1\}$, the homotopy is
\[
 H'(\chi^e\otimes u)
       =(-1)^{3e}\chi^e
                                      \otimes H_{\mathrm{rel}}u.
\]
For two successive contractions, the composite homotopy is
$H_A+i_AH_Bp_A$.
Expansion of this expression proves the contraction identity;
the mixed terms in its square vanish by the side conditions.
This specifies the final homotopy $H$ on every polynomial input.

\begin{theorem}\label{thm:vol2-absolute-contraction}
There is a normalized contraction over $R$
\[
 (\mathcal B_\hbar,Q)
    \underset{i}{\overset{\Pi}{\rightleftarrows}}(Z,0),
 \qquad QH+HQ=1-i\Pi,
 \qquad \Pi i=1,\quad H^2=Hi=\Pi H=0.
\]
It is polynomial on each input and remains a contraction
after every commutative coefficient base change $R\to R'$.
For $F=\sum_I\gamma^IF_I$ in normal order,
\begin{equation}\label{eq:vol2-absolute-projection}
 \Pi(F)=\rho_\hbar(P_0F_\varnothing)
            +\chi\rho_\hbar(P_0F_{123}),
 \qquad
 i(u+\chi v)=i_{\mathrm{rel}}u+\chi i_{\mathrm{rel}}v.
\end{equation}
\end{theorem}

\begin{proof}
The three contractions just constructed have precisely these
inclusion and projection maps. Composition proves all displayed
identities. The Casimir operators use finite-dimensional
coefficient representations and only nonzero complex integer
denominators. The relative operator $T^{-1}$ terminates on
each polynomial, as do the finite weight splittings.
Every constituent is $R$-linear with polynomial output.
Tensoring the identities with $R'$ therefore preserves them,
without a flatness assumption on $R'$.
The resulting homotopy also proves that the projection kernel
is contractible, rather than merely giving a list of proposed
cohomology representatives.
\end{proof}

Let $E=W_\hbar(s,t)^{\mathrm{even}}$ and
$M=W_\hbar(s,t)^{\mathrm{odd}}$, where parity means total
polynomial degree modulo two. The cohomology groups are
\begin{equation}\label{eq:vol2-absolute-cohomology-degrees}
 H^{-1}=M\theta,
 \qquad H^0=E,
 \qquad H^2=\chi M\theta,
 \qquad H^3=\chi E,
\end{equation}
and all other degrees vanish. The odd module and its singular
fibre from the relative calculation remain present.
The absolute operation adds the two positive degrees.
This is a polynomial statement. Extending the contraction to
a completion requires its own continuity and convergence
conditions.

\section{Products and the failure of a strict projection}
\label{sec:vol2-absolute-products}

Give $Z=\Lambda(\chi)\otimes\mathcal H_\hbar$ the graded
tensor product. Thus
$u\chi=(-1)^{|u|}\chi u$ and $\chi^2=0$.
The product on $\mathcal H_\hbar$ is its Weyl--exterior
product from the relative restriction.

\begin{proposition}\label{prop:vol2-absolute-ring}
Under the contraction, the cohomology algebra is
\[
 H^*(\mathcal B_\hbar,Q)
       \cong\Lambda(\chi)\otimes\mathcal H_\hbar.
\]
This is an isomorphism of cohomology rings. The linear
projection $\Pi$ is not an algebra map on the full complex.
\end{proposition}

\begin{proof}
Every relative class has an invariant closed representative
of antighost degree zero or one. Products of representatives
with no top ghost have the relative cohomology product.
For a product with one top ghost, normal ordering preserves
the term with three positive ghosts, with exactly the graded
tensor sign. All contraction terms have positive-ghost
degree two and are killed by $\Pi$.
The relative projection of the remaining coefficient is
its relative product.

Products with two top ghosts and at most two antighosts
are zero. Six positive ghosts cannot survive among three
exterior generators unless at least three are removed by
contractions. Only two antighosts are available.
This proves all products in the asserted ring.
Since the representatives are closed and $\Pi$ is the
projection of a contraction, these projected products compute
their actual cohomology classes.
\end{proof}

The graded commutation with $\chi$ is a cohomology identity.
Its cochain defect has an explicit primitive:
\begin{equation}\label{eq:vol2-absolute-top-commutators}
 \chi\star\xi+\xi\star\chi
     =\hbar Q\left(\sum_i x_i\gamma^i\right),
 \qquad
 \chi\star\eta+\eta\star\chi
     =\hbar Q\left(\sum_i y_i\gamma^i\right).
\end{equation}
For the first equation, normal ordering gives the left side as
\[
 \hbar\left(x_1\gamma^2\gamma^3
       -x_2\gamma^1\gamma^3+x_3\gamma^1\gamma^2\right).
\]
The coefficient-action part of $Q(\sum x_i\gamma^i)$
gives twice the expression in parentheses, and the scalar
Chevalley part subtracts it once. The antighost differential
does not enter. This proves the first identity, including its
sign; replacing $x$ by $y$ proves the second.

\begin{proposition}\label{prop:vol2-absolute-no-strict-projection}
There is no graded unital $R$-algebra map
$\mathcal B_\hbar\to Z$.
\end{proposition}

\begin{proof}
Equation~\eqref{eq:vol2-absolute-cohomology-degrees} shows
$Z^1=0$. Such a map would send every $\gamma^i$ to zero.
The Clifford relation
$\psi_i\gamma^i+\gamma^i\psi_i=\hbar$ would then send
$\hbar$ to zero. A unital $R$-linear map must preserve it,
and $\hbar\ne0$ in the $R$-free target. This is a contradiction.
\end{proof}

The assertion fixes the direction of the map and the coefficient
ring. At the classical specialization $\hbar=0$, its Clifford
obstruction disappears. It also says nothing about a sequence
of quasi-isomorphisms whose arrows may point in either direction.
A dg algebra is \emph{formal} if such a sequence connects it
to its cohomology ring with zero differential. Determining
formality of this absolute algebra requires an invariant of
that comparison problem.

\section{The higher model}
\label{sec:vol2-absolute-higher-model}

The contraction provides the operation data beyond the ring.
Let $s$ denote suspension of degree $-1$, distinct from the
Cartan coordinate in $W_\hbar(s,t)$, and put
\[
 V=s\mathcal B_\hbar,\qquad W=sZ,
 \qquad d_V=-sQs^{-1},\qquad h_V=-sHs^{-1},
 \qquad i_s=sis^{-1},\quad p_s=s\Pi s^{-1}.
\]
For a graded module $U$, use the tensor coalgebra
$T^c(U)=\bigoplus_{n\geq1}U^{\otimes n}$ with the
deconcatenation coproduct. Tensor products of maps use the
Koszul sign. On $T^c(V)$ let $d_T$ be the tensor differential
and let the multiplication coderivation be determined by
\begin{equation}\label{eq:vol2-absolute-bar-multiplication}
 b_2(sa,sb)=(-1)^{|a|}s(a\star b).
\end{equation}
The derivation property of $Q$ and associativity give
$(d_T+b_2)^2=0$.
This uses the underlying algebra with its unit forgotten;
it assumes no scalar augmentation of a Weyl algebra.

On tensor length $n$ extend the contraction by
\begin{equation}\label{eq:vol2-absolute-tensor-homotopy}
 i_T=i_s^{\otimes n},\quad p_T=p_s^{\otimes n},\qquad
 \mathbf H=
 \sum_{j=1}^n(i_sp_s)^{\otimes(j-1)}
                     \otimes h_V\otimes1^{\otimes(n-j)}.
\end{equation}
The identity $d_Vh_V+h_Vd_V=1-i_sp_s$ telescopes to
\[
 d_T\mathbf H+\mathbf H d_T=1-i_Tp_T,
 \qquad p_Ti_T=1,
 \qquad \mathbf H^2=\mathbf H i_T=p_T\mathbf H=0.
\]
For instance, in a mixed term of $\mathbf H^2$ the earlier
homotopy is followed by $i_sp_s$ or the later one is preceded
by it. The factors $h_Vi_s$ and $p_sh_V$ kill these terms;
the equal-position term contains $h_V^2$.

\begin{theorem}\label{thm:vol2-absolute-higher-model}
The finite-at-each-length formulas
\begin{align}\label{eq:vol2-absolute-transfer}
 b_Z&=p_Tb_2(1+\mathbf H b_2)^{-1}i_T,
       \notag\\
 \mathbf I&=(1+\mathbf H b_2)^{-1}i_T,
 &\mathbf P&=p_T(1+b_2\mathbf H)^{-1}
\end{align}
give a minimal $A_\infty$ structure on $Z$ and
$A_\infty$ comparison maps with $(\mathcal B_\hbar,Q)$.
The binary product is the ring in
Proposition~\ref{prop:vol2-absolute-ring}.
All components are polynomial over $R$ on polynomial inputs
and persist under coefficient base change.
The construction has a strictly unital normalized presentation.
\end{theorem}

\begin{proof}
Each occurrence of $b_2$ reduces tensor length by one.
Thus every inverse in \eqref{eq:vol2-absolute-transfer}
is a finite geometric sum on a tensor of fixed length $n$,
with at most $n-1$ occurrences of $b_2$.
Apply the perturbation identity of
Lemma~\ref{lem:vol2-filtered-perturbation} to the tensor
contraction. Its finite sums give
\[
 b_Z^2=0,
 \quad (d_T+b_2)\mathbf I=\mathbf I b_Z,
 \quad \mathbf P(d_T+b_2)=b_Z\mathbf P,
 \quad \mathbf P\mathbf I=1.
\]
The comparison homotopy is
$\mathbf H'=(1+\mathbf H b_2)^{-1}\mathbf H$.

We also need the coalgebra structures, not just these complex
identities. Cutting a tensor in
\eqref{eq:vol2-absolute-tensor-homotopy} on either side of
its homotopy factor gives
\begin{equation}\label{eq:vol2-absolute-homotopy-coproduct}
 \Delta\mathbf H
   =\bigl(\mathbf H\otimes1+i_Tp_T\otimes\mathbf H\bigr)\Delta.
\end{equation}
The maps $i_T,p_T$ are coalgebra maps, and $b_2$ is a
coderivation. Insert these three identities into the finite
series. Induction on the number of multiplications gives
\[
 \Delta\mathbf I=(\mathbf I\otimes\mathbf I)\Delta,
 \qquad
 \Delta\mathbf P=(\mathbf P\otimes\mathbf P)\Delta,
 \qquad
 \Delta b_Z=(b_Z\otimes1+1\otimes b_Z)\Delta.
\]
To describe the induction explicitly, cut each term at one
of its surviving tensor gaps. The multiplication is on the
left or right of that gap, so the two choices give the two
tensor factors of the preceding identities. A homotopy crossing
the cut splits as the two terms of
\eqref{eq:vol2-absolute-homotopy-coproduct}.
Terms with $\mathbf H i_T$, $p_T\mathbf H$ or
$\mathbf H^2$ vanish. The remaining left and right strings
are exactly the shorter geometric series. Their tensor signs
are inherited from the displayed coproduct identity.
This proves the induction at every length.

There is no unary component of $b_Z$, since the cohomology
differential is zero. The binary component is
$p_sb_2(i_s\otimes i_s)$, which is the computed cohomology
product. Projection to length one gives all higher components.
Each is a finite sum of products and the polynomial contraction,
so it has the asserted coefficient and specialization properties.

Finally $i(1)=1$, $\Pi(1)=1$ and $H(1)=0$.
In a higher tree with a unit input, the first multiplication
meeting that input removes the unit. An adjacent homotopy
then gives $Hi=0$ or $H^2=0$, or the final projection gives
$\Pi H=0$. This is the normalized strict-unit construction:
the binary unit laws remain, while higher operations and higher
comparison components with a unit input vanish.
It removes the tensors degenerate by a unit insertion, without
introducing an algebra augmentation.
\end{proof}

The formulas specify every operation once the relative weight
splittings are chosen. They do not assert that the higher
operations vanish. For the next calculation, keep the relative
section and homotopy fixed in
Section~\ref{sec:vol2-absolute-contraction}.

\subsection{A ternary operation from two trees}

The chosen full homotopy has
\begin{equation}\label{eq:vol2-absolute-tree-homotopies}
 H(\xi\star\eta)=r,
 \qquad
 H(\eta\star\chi)=\hbar\sum_i y_i\gamma^i.
\end{equation}
The first equality is the relative choice.
For the second, write
\[
 \eta\star\chi=-\chi\eta+\hbar\omega_y,
 \qquad
 \omega_y=y_1\gamma^2\gamma^3
          -y_2\gamma^1\gamma^3+y_3\gamma^1\gamma^2.
\]
The first term is a chosen cohomology representative.
On the coefficient spin-one summand,
$G=-1/2$. Direct rotation gives
$K\omega_y=-2\sum_i y_i\gamma^i$ and
$d_{\mathrm{CE}}(\sum_i y_i\gamma^i)=\omega_y$.
Thus the normalized Casimir homotopy sends $\omega_y$ to
$\sum_i y_i\gamma^i$, proving the second equality.

\begin{proposition}\label{prop:vol2-absolute-ternary}
For this contraction, the suspended ternary coefficient is
\begin{equation}\label{eq:vol2-absolute-ternary}
 (b_Z)_3\bigl(s(s\theta),s(t\theta),s\chi\bigr)
     =s\left(\hbar^2st+\frac12\hbar^3\right).
\end{equation}
Inside the parentheses, $s,t$ are the Cartan Weyl symbols.
The outer $s$ is suspension.
\end{proposition}

\begin{proof}
The three suspended inputs have even parity.
The length-three term of \eqref{eq:vol2-absolute-transfer}
is $-p_Tb_2\mathbf H b_2i_T$.
Using \eqref{eq:vol2-absolute-tree-homotopies} and the minus
sign in $h_V=-sHs^{-1}$, its two unsuspended contributions
have a plus sign and are
\[
 \Pi(r\star\chi),\qquad
 \Pi\left(\xi\star\hbar\sum_i y_i\gamma^i\right).
\]
The scalar part of
$\psi_1\psi_2\psi_3\gamma^1\gamma^2\gamma^3$
is $-\hbar^3$. All other terms have one or two positive
ghosts, or have three positive ghosts and the relative factor
$r$, whose projection is zero. Hence
\[
 \Pi(r\star\chi)=-\hbar^3.
\]
In the second tree, only the contractions
$\psi_i\gamma^j\mapsto\hbar\delta_i^j$
contribute to positive-ghost degree zero.
The Weyl product gives
$x_i\star y_i=x_iy_i+\hbar/2$.
The invariant sum $b=\sum_i x_iy_i$ restricts to $st$,
so this tree has projection
\[
 \hbar^2st+\frac32\hbar^3.
\]
Their sum proves the formula.
\end{proof}

Both the complete ghost contraction and the Weyl scalar term
are necessary. The coefficient is a property of this specified
minimal model. To use it as a formality obstruction, one would
have to show that its Hochschild class cannot be removed by
an $A_\infty$ change of coordinates. A single displayed value
does not prove that assertion.
The polynomial-subalgebra section in
Proposition~\ref{prop:vol2-relative-polynomial-section}
already shows how particular relative higher values can change
without changing the underlying quasi-isomorphism type.

\section{The central matrix direction}
\label{sec:vol2-absolute-central-direction}

The traceless calculation uses $\mathfrak{sl}_2$.
For $\mathfrak{gl}_2$, split off the scalar matrix direction
using the trace pairing. Normalize its coordinates $x_0,y_0$
so that $[x_0,y_0]=\hbar$.
Its moment function and adjoint action are zero.
The corresponding ghosts therefore have zero differential,
but their product is
\begin{equation}\label{eq:vol2-absolute-central-clifford}
 \psi_0\gamma^0+\gamma^0\psi_0=\hbar,
 \qquad (\gamma^0)^2=\psi_0^2=0,
 \qquad |\gamma^0|=1,\quad |\psi_0|=-1.
\end{equation}
In this split normalization, the full complex is
\[
 (\mathcal B_\hbar,Q)\otimes_R W_\hbar(x_0,y_0)
       \otimes_R\bigl(\mathrm{Cl}_\hbar(\gamma^0,\psi_0),0\bigr).
\]
Both additional factors are $R$-free with zero differential.
Tensoring the constructed contraction computes its cohomology
ring as the same two factors tensored with $Z$.

At $\hbar=0$ the central ghost product specializes to the
exterior product. Over $R$ it retains
\eqref{eq:vol2-absolute-central-clifford}.
After inverting $\hbar$, the central Clifford algebra acts as
the full endomorphism algebra of the two-dimensional graded
space $\Lambda(\gamma^0)$.
This may be used for a Morita comparison of module categories;
it does not identify the algebra with an exterior algebra by
a dg-algebra quasi-isomorphism.

\section{Absolute formality and the Hamiltonian comparison}
\label{sec:vol2-absolute-remaining-comparison}

The contraction determines the absolute complex, its ring
and one complete higher model over $\mathbb C[\hbar]$.
The formality question asks for a zigzag to the displayed
cohomology ring, or an obstruction invariant under every
allowed change of higher model. The first possible higher
product must be treated as a Hochschild class, including its
indeterminacy under such changes. The strict-map obstruction
and the chosen ternary value have their narrower conclusions.

The Hamiltonian comparison requires additional structure.
Its first component must agree with the global ordinary trace
map, remain strict on the scalar quadratic sector, and retain
the cubic primitive of
Section~\ref{sec:vol2-derived-action-bfv}.
The higher equations must be solved on this full absolute
recipient, with the negative relation classes, positive gauge
classes and cyclic operations present.
The matrix parameter $\hbar$ is independent of the
coefficient-plane parameter $\varepsilon$.
A linear contraction does not supply those equivariances.

Nor does the rank-two calculation give a higher-rank absolute
contraction or a field-theoretic realization.
Those comparisons must specify their own relation complexes,
supported operations and coefficient completions.
The explicit recipient constructed here makes their required
maps precise, including the classes that an ordinary
degree-zero reduction would lose.
